\section{Results}
\label{sec:results}

\subsection{Fitted observation-model spread}
\label{sec:model-spread}

Figure~\ref{fig:spectrum-models} shows the first real spectrum and all fitted observation candidates. For $x=0.226$, the fitted fractional-power and Chu candidates span 181.086--187.500~meV, a full range of 6.414~meV. The $p=0.5$ and Chu candidates reach the upper edge-search boundary and are flagged as boundary-limited.

For $x=0.310$, the fitted candidates span 289.670--296.500~meV, a full range of 6.830~meV. The Chu candidate is boundary-limited. The model-family spread is therefore similar for the two compositions despite their different nominal gaps and film thicknesses. No single fitted edge is selected.

\begin{figure*}[t]
  \centering
  \IfFileExists{figures/figure1_spectrum_models.pdf}{%
    \includegraphics[width=0.92\textwidth]{figures/figure1_spectrum_models.pdf}%
  }{\fbox{\parbox[c][55mm][c]{0.88\textwidth}{\centering Figure 1 PDF export pending from frozen SVG}}}
  \caption{Digitized Moazzami 2005 Fig.~6a IRSE spectrum at $x=0.226$ and 300~K with every fitted observation-model candidate. Boundary-limited candidates are retained as diagnostics. No fitted curve is selected as the latent material gap.}
  \label{fig:spectrum-models}
\end{figure*}

\subsection{Digitization sensitivity does not explain the model spread}
\label{sec:digitization-result}

Across both spectra, the maximum coherent coordinate-induced shift of any fitted-model edge is 0.891~meV. This is substantially smaller than the 6.414--6.830~meV fitted-model spans. The observation-model dependence therefore survives the declared digitization uncertainty.

Every fixed threshold through 4000~cm$^{-1}$ remains below the 5~meV coordinate-shift gate. The 5000~cm$^{-1}$ crossing for $x=0.310$ shifts by 5.694~meV and is excluded from precision plots and threshold-based claims.

\subsection{Observation definition changes the nominal material comparator}
\label{sec:ranking-result}

Figure~\ref{fig:edge-candidates} displays all extracted edges by observation definition. For every fitted fractional-power and Chu candidate, the published Seiler relation is nominally closest for both specimens. Fixed absorption crossings produce a different result.

For $x=0.226$, published Seiler is nominally closest at 400 and 600~cm$^{-1}$, Laurenti is closest near 800--1000~cm$^{-1}$, and the provisional Hansen--Pade relation is closest at higher crossings. For $x=0.310$, published Seiler remains nominally closest through 1000~cm$^{-1}$, the provisional relation is closest at 2000~cm$^{-1}$, and Laurenti is closest at 3000~cm$^{-1}$ and above.

Thus, a material-model ranking can be reversed by changing only the declared observation definition. This does not imply that high-threshold crossings better estimate the latent gap. It demonstrates that the ranking is conditional on what was called the edge.

\begin{figure*}[t]
  \centering
  \IfFileExists{figures/figure2_edge_candidates.pdf}{%
    \includegraphics[width=0.98\textwidth]{figures/figure2_edge_candidates.pdf}%
  }{\fbox{\parbox[c][58mm][c]{0.94\textwidth}{\centering Figure 2 PDF export pending from frozen SVG}}}
  \caption{Extracted edge versus observation definition for $x=0.226$ and $x=0.310$. The upper group contains fitted fractional-power and Chu candidates; the lower group contains fixed absorption crossings. Open markers denote boundary-limited fits.}
  \label{fig:edge-candidates}
\end{figure*}

\subsection{The Seiler--Hansen fitted-edge ordering is not identified}
\label{sec:seiler-result}

At $x=0.226$ and 300~K, the Hansen and published Seiler predictions are 190.366 and 190.112~meV, respectively, a separation of 0.255~meV. At $x=0.310$ and 300~K, the predictions are 304.235 and 304.059~meV, a separation of 0.177~meV.

Published Seiler is therefore nominally closer to every fitted-model edge, but the margin over Hansen is sub-meV. The source does not report specimen-level composition uncertainty, and the two spectra are not independent cross-laboratory records. The data do not identify a strict Seiler-over-Hansen material-law ranking. Figure~\ref{fig:residual-intervals} places that sub-meV comparator separation in the context of the wider observation-definition intervals.

\begin{figure*}[t]
  \centering
  \IfFileExists{figures/figure3_material_residual_envelopes.pdf}{%
    \includegraphics[width=0.92\textwidth]{figures/figure3_material_residual_envelopes.pdf}%
  }{\fbox{\parbox[c][55mm][c]{0.88\textwidth}{\centering Figure 3 PDF export pending from frozen SVG}}}
  \caption{Residual intervals relative to Hansen, published Seiler, Laurenti, and provisional Hansen--Pade predictions. Narrow dark intervals show the fitted-model envelope; wider gray intervals show stable 400--4000~cm$^{-1}$ threshold definitions. These intervals represent observation-definition sensitivity, not statistical confidence intervals.}
  \label{fig:residual-intervals}
\end{figure*}

\subsection{Audit-grade acquisition target}
\label{sec:design-result}

The identifiability structure is summarized in Fig.~\ref{fig:identifiability}. The paired design in Fig.~\ref{fig:acquisition-design} separates the common latent gap from observation-class, carrier, and vacancy terms by construction. Eight specimen states are repeated at low and room temperature, and both modalities are applied to every specimen. Independent composition metrology, Hall state, a vacancy-sensitive proxy, raw absorption data, extraction covariance, and processing history are required.

The design does not assume that the five-term local model is physically complete. It supplies the minimum redundancy needed to test that model under the declared gates.

\begin{figure*}[t]
  \centering
  \IfFileExists{figures/figure4_identifiability.pdf}{%
    \includegraphics[width=0.98\textwidth]{figures/figure4_identifiability.pdf}%
  }{\fbox{\parbox[c][58mm][c]{0.94\textwidth}{\centering Figure 4 PDF export pending from frozen SVG}}}
  \caption{Identifiability diagram separating the latent material gap from method, carrier, vacancy, and measurement terms. Same-specimen magneto-optical and absorption observations remove the common latent term in the paired difference. The decomposition does not authorize a universal correction.}
  \label{fig:identifiability}
\end{figure*}

\begin{figure*}[t]
  \centering
  \IfFileExists{figures/figure5_paired_acquisition_design.pdf}{%
    \includegraphics[width=0.98\textwidth]{figures/figure5_paired_acquisition_design.pdf}%
  }{\fbox{\parbox[c][58mm][c]{0.94\textwidth}{\centering Figure 5 PDF export pending from frozen SVG}}}
  \caption{Audit-grade paired $2\times2\times2$ acquisition design. Eight composition--carrier--vacancy specimen states receive paired magneto-optical and absorption measurements in low- and room-temperature blocks, producing 32 observations. Numerical diagnostics are read directly from the validated design oracle.}
  \label{fig:acquisition-design}
\end{figure*}

\begin{table*}[t]
\centering
\small
\setlength{\tabcolsep}{3pt}
\caption{Specimen and measurement provenance for the two real-spectrum records. Blank composition uncertainty and unknown carrier state are retained as missing metadata rather than inferred. Full source and input hashes remain in the frozen CSV record.}
\label{tab:specimen-provenance}
\begin{tabularx}{\textwidth}{@{}rrrrp{0.13\textwidth}Yp{0.19\textwidth}@{}}
\toprule
$x$ & $T$ (K) & Thickness ($\mu$m) & Points & $\sigma_x$ & Carrier state & Source \\
\midrule
0.226 & 300 & 15.40 & 125 & not reported & unknown; density not reported & Moazzami 2005, Fig.~6a \\
0.310 & 300 & 4.95 & 115 & not reported & unknown; density not reported & Moazzami 2005, Fig.~6b \\
\bottomrule
\end{tabularx}
\end{table*}

\begin{table*}[t]
\centering
\caption{Observation-candidate definitions and source-supported domains. Operational thresholds are not identified with the latent material gap.}
\label{tab:candidate-definitions}
\begin{tabular}{p{0.22\textwidth}p{0.39\textwidth}p{0.29\textwidth}}
\toprule
Candidate & Definition & Admissible domain \\
\midrule
fractional power, fitted $p$ & $\alpha=A(E-E_g)^p/E$; $p$ fitted & declared fit window \\
fractional power, fixed $p$ & $\alpha=A(E-E_g)^p/E$; $p$ fixed & $p=0.5$, $0.7$, source-panel $p$, and $1.0$ \\
Chu 1994 Kane region & $\alpha=\alpha_g\exp[\sqrt{\beta(x,T)(E-E_g)}]$ & $0.170\leq x\leq0.443$; $77\leq T\leq300$~K \\
fixed absorption threshold & First interpolated $\alpha$ crossing & 400--5000~cm$^{-1}$; 5000 flagged when coordinate-sensitive \\
\bottomrule
\end{tabular}
\end{table*}

\begin{landscape}
\begingroup
\scriptsize
\renewcommand{\arraystretch}{0.95}
\setlength{\tabcolsep}{3pt}
\setlength{\LTleft}{0pt}
\setlength{\LTright}{0pt}
\begin{longtable}{@{}rp{1.70in}p{0.65in}rrrrp{1.50in}@{}}
\caption{Complete extracted-edge ensemble. Coordinate shift is the maximum absolute shift over the declared coherent digitization perturbations. Boundary-limited candidates are retained but are not treated as identified edges.}\label{tab:edge-ensemble}\\
\toprule
$x$ & Candidate & Type & Edge (meV) & Bound? & Shift (meV) & Lead (meV) & Closest comparator \\
\midrule
\endfirsthead
\multicolumn{8}{c}{\tablename\ \thetable{} -- continued}\\
\toprule
$x$ & Candidate & Type & Edge (meV) & Bound? & Shift (meV) & Lead (meV) & Closest comparator \\
\midrule
\endhead
\midrule
\multicolumn{8}{r}{Continued on next page}\\
\endfoot
\bottomrule
\endlastfoot
0.226 & fractional power, fitted $p$ & fit & 181.711 & no & 0.623 & 0.255 & Seiler 1990 \\
0.226 & fractional $p=0.5$ & fit & 187.500 & yes & 0.526 & 0.255 & Seiler 1990 \\
0.226 & fractional $p=0.7$ & fit & 186.180 & no & 0.643 & 0.255 & Seiler 1990 \\
0.226 & fractional $p=0.872$ & fit & 183.691 & no & 0.774 & 0.255 & Seiler 1990 \\
0.226 & fractional $p=1$ & fit & 181.086 & no & 0.891 & 0.255 & Seiler 1990 \\
0.226 & Chu 1994 Kane region & fit & 187.500 & yes & 0.526 & 0.255 & Seiler 1990 \\
0.226 & $\alpha=400$~cm$^{-1}$ & threshold & 183.124 & no & 1.181 & 0.255 & Seiler 1990 \\
0.226 & $\alpha=600$~cm$^{-1}$ & threshold & 187.705 & no & 1.131 & 0.255 & Seiler 1990 \\
0.226 & $\alpha=800$~cm$^{-1}$ & threshold & 191.281 & no & 1.287 & 0.651 & Laurenti \\
0.226 & $\alpha=1000$~cm$^{-1}$ & threshold & 194.090 & no & 1.283 & 0.651 & Laurenti \\
0.226 & $\alpha=2000$~cm$^{-1}$ & threshold & 207.154 & no & 1.880 & 8.340 & Provisional Hansen--Pade \\
0.226 & $\alpha=3000$~cm$^{-1}$ & threshold & 223.091 & no & 2.965 & 8.340 & Provisional Hansen--Pade \\
0.226 & $\alpha=4000$~cm$^{-1}$ & threshold & 244.961 & no & 3.184 & 8.340 & Provisional Hansen--Pade \\
0.226 & $\alpha=5000$~cm$^{-1}$ & threshold & 270.951 & no & 3.954 & 8.340 & Provisional Hansen--Pade \\
\addlinespace
0.310 & fractional power, fitted $p$ & fit & 289.670 & no & 0.553 & 0.177 & Seiler 1990 \\
0.310 & fractional $p=0.5$ & fit & 296.348 & no & 0.546 & 0.177 & Seiler 1990 \\
0.310 & fractional $p=0.7$ & fit & 295.248 & no & 0.579 & 0.177 & Seiler 1990 \\
0.310 & fractional $p=0.849$ & fit & 293.784 & no & 0.626 & 0.177 & Seiler 1990 \\
0.310 & fractional $p=1$ & fit & 291.763 & no & 0.679 & 0.177 & Seiler 1990 \\
0.310 & Chu 1994 Kane region & fit & 296.500 & yes & 0.526 & 0.177 & Seiler 1990 \\
0.310 & $\alpha=400$~cm$^{-1}$ & threshold & 292.225 & no & 1.122 & 0.177 & Seiler 1990 \\
0.310 & $\alpha=600$~cm$^{-1}$ & threshold & 296.482 & no & 1.268 & 0.177 & Seiler 1990 \\
0.310 & $\alpha=800$~cm$^{-1}$ & threshold & 299.858 & no & 1.292 & 0.177 & Seiler 1990 \\
0.310 & $\alpha=1000$~cm$^{-1}$ & threshold & 302.322 & no & 1.160 & 0.177 & Seiler 1990 \\
0.310 & $\alpha=2000$~cm$^{-1}$ & threshold & 312.722 & no & 1.938 & 1.786 & Provisional Hansen--Pade \\
0.310 & $\alpha=3000$~cm$^{-1}$ & threshold & 322.693 & no & 1.693 & 6.289 & Laurenti \\
0.310 & $\alpha=4000$~cm$^{-1}$ & threshold & 333.492 & no & 3.643 & 6.289 & Laurenti \\
0.310 & $\alpha=5000$~cm$^{-1}$ & threshold & 351.695 & no & 5.694 & 6.289 & Laurenti \\
\end{longtable}
\endgroup
\end{landscape}

\begin{table*}[t]
\centering
\small
\setlength{\tabcolsep}{4pt}
\caption{Material-gap predictions and fitted observation-model residual intervals. Residual is extracted edge minus model prediction. Strict material-law ranking is not authorized for either specimen.}
\label{tab:model-comparison}
\begin{tabularx}{\textwidth}{@{}rYrp{0.27\textwidth}c@{}}
\toprule
$x$ & Comparator & Prediction (meV) & Fitted residual interval (meV) & Strict ranking? \\
\midrule
0.226 & Hansen & 190.366 & [-9.281, -2.866] & no \\
0.226 & Seiler 1990 & 190.112 & [-9.026, -2.612] & no \\
0.226 & Laurenti & 191.018 & [-9.932, -3.518] & no \\
0.226 & Provisional Hansen--Pade & 199.357 & [-18.272, -11.857] & no \\
\addlinespace
0.310 & Hansen & 304.235 & [-14.565, -7.735] & no \\
0.310 & Seiler 1990 & 304.059 & [-14.389, -7.559] & no \\
0.310 & Laurenti & 316.759 & [-27.089, -20.259] & no \\
0.310 & Provisional Hansen--Pade & 310.470 & [-20.800, -13.970] & no \\
\bottomrule
\end{tabularx}
\end{table*}

\begin{table*}[t]
\centering
\caption{Claim and evidence boundary for the present manuscript.}
\label{tab:claim-boundaries}
\begin{tabular}{p{0.40\textwidth}p{0.15\textwidth}p{0.35\textwidth}}
\toprule
Claim & Status & Evidence boundary \\
\midrule
fractional-model edge sensitivity is approximately 6--7 meV & authorized & two digitized 300~K IRSE spectra from one source study \\
published Seiler is nominally closest for fitted-model edges & descriptive only & 0.18--0.25 meV advantage over Hansen; composition uncertainty missing \\
fixed-threshold definition can change the nominal comparator & authorized through 4000~cm$^{-1}$ & threshold is operational and is not the latent material gap \\
one corrected or production edge exists & not authorized & complete candidate ensemble must be reported \\
\bottomrule
\end{tabular}
\end{table*}

